\documentclass[11pt]{article}

\usepackage[letterpaper,margin=1in]{geometry}
\usepackage[T1]{fontenc}
\usepackage{lmodern}
\usepackage{amsmath}
\usepackage{amssymb}
\usepackage{amsthm}
\usepackage{enumitem}
\usepackage{natbib}
\usepackage[hidelinks]{hyperref}

\newtheorem{proposition}{Proposition}
\newtheorem{theorem}{Theorem}
\newtheorem{lemma}{Lemma}
\newtheorem{corollary}{Corollary}
\theoremstyle{definition}
\newtheorem{assumption}{Assumption}
\newtheorem{definition}{Definition}
\theoremstyle{remark}
\newtheorem{remark}{Remark}

\setlist[enumerate]{itemsep=1pt,topsep=3pt}

\newcommand{\cN}{\mathcal N}
\newcommand{\cK}{\mathcal K}
\newcommand{\cG}{\mathcal G}
\newcommand{\cP}{\mathcal P}
\newcommand{\cA}{\mathcal A}
\newcommand{\IO}{\mathrm{IO}}

\begin{document}

\section{Model}
\label{sec:model}

We develop a static general equilibrium model of multinational production in which
the strategic agent is the ownership group. A group can operate plants in several
countries, chooses which of its plants serve each market, and, when it sets
quantities, takes into account that its own plants compete with one another. Markups
therefore depend on the group's total share of the market rather than on the share of
any single plant. Profits are paid to the countries that own each group, and ownership
is recorded group by group rather than country by country. Proofs are in the
Appendix.

\subsection{Model Set-Up}
\label{sec:setup}

\paragraph{Agents.} There is a finite set of countries $\cN=\{1,\dots,N\}$ and a
finite set of sectors $\cK=\{1,\dots,K\}$. Each country $n$ is inhabited by a
representative household that supplies $L_n$ units of labor inelastically and owns
shares in firms. Production takes place in plants, and every plant belongs to exactly
one group $g\in\cG$, which we identify with its global ultimate parent. Each group has
a headquarters country $h(g)$ and operates in a single sector
$k(g)$.\footnote{This restriction is without loss for behavior. Marginal costs are
constant and fixed costs are specific to each plant--market pair, so a group's
decisions separate across markets (Section~\ref{sec:market}). A group active in
several sectors is equivalent to several single-sector groups with the same owners;
what the restriction rules out is a capability shared across sectors.} Groups are of
two types. A \emph{parent group} $g\in\cG^P$ is headquartered in a country of a set
$\cN^P\subseteq\cN$ and has one potential plant in every country of $\cN^P$, its
headquarters country included. A \emph{stand-alone firm} $g\in\cG^S$ has a single
potential plant, located in $h(g)$; countries outside $\cN^P$ host only stand-alone
firms. We index plants by $i$, write $g(i)$ and $\ell(i)$ for the group and the
location of plant $i$ and $k(i)=k(g(i))$ for its sector, and let $\cP_g$ be the set of
potential plants of group $g$. Each plant produces one differentiated variety and can
sell it in any country. A \emph{market} is a destination--sector pair $(n,k)$.
Whether a group is multinational is an equilibrium outcome: a group is multinational if
its active plants are located in more than one country, an active plant located outside
its group's headquarters country is a foreign affiliate, and the home plant of a
multinational group is a domestic multinational. A stand-alone firm is never
multinational.

\paragraph{Representative household.} The household in country $n$ has nested CES
preferences,
\begin{equation}
  U_n=\Bigl[\sum_{k\in\cK}\beta_k^{1/\eta}\,C_{nk}^{(\eta-1)/\eta}\Bigr]^{\eta/(\eta-1)},
  \qquad
  C_{nk}=\Bigl[\sum_{i\in\cA_{nk}}q_{i,n}^{(\sigma_k-1)/\sigma_k}\Bigr]^{\sigma_k/(\sigma_k-1)},
  \label{eq:utility}
\end{equation}
where $\cA_{nk}$ is the set of plants active in market $(n,k)$, $\beta_k>0$ with
$\sum_k\beta_k=1$, and $\sigma_k>\eta\ge1$: varieties within a sector are closer
substitutes than goods in different sectors. The case $\eta=1$ is the Cobb--Douglas
limit. The associated price indices are
\begin{equation}
  P_{nk}=\Bigl[\sum_{i\in\cA_{nk}}p_{i,n}^{1-\sigma_k}\Bigr]^{1/(1-\sigma_k)},
  \qquad
  P_n=\Bigl[\sum_{k}\beta_k P_{nk}^{1-\eta}\Bigr]^{1/(1-\eta)} .
  \label{eq:priceindex}
\end{equation}
The household spends its income $X_n$, allocates the share
$\beta_kP_{nk}^{1-\eta}/P_n^{1-\eta}$ of it to sector $k$, and its welfare is real
income $X_n/P_n$. Firms buy intermediate inputs as the same sectoral composites (below), so
all purchases of sector-$k$ varieties in country $n$, by households and by firms, are
aggregated by $C_{nk}$. Writing $q_{i,n}$ for the total quantity of plant $i$'s
variety sold in $n$ and $E_{nk}$ for total spending on sector $k$ in $n$, demand is
\begin{equation}
  q_{i,n}=p_{i,n}^{-\sigma_k}\,P_{nk}^{\sigma_k-1}E_{nk}.
  \label{eq:demand}
\end{equation}
Market sizes $E_{nk}$ are determined in general equilibrium (Section~\ref{sec:ge}).%
\footnote{Intermediates are combined across sectors with Cobb--Douglas weights, so
when they are used ($\nu_k>0$ below) we set $\eta=1$: final and intermediate demand
then share the same upper-tier elasticity, and the markup in \eqref{eq:markupC} is the
one firms face. The case $\eta>1$ is used only without intermediates.}

\paragraph{Production technology.} Labor is the only primary factor. A plant in sector
$k$ located in $\ell$ combines labor with a bundle of intermediate inputs. Writing
$h_i\equiv h(g(i))$ for the headquarters country of plant $i$, its unit cost of
production is
\begin{equation}
  a_i=\frac{1}{\varphi_i}\,w_\ell^{\,1-\nu_k}\,\bigl(P^{\IO}_{\ell k}\bigr)^{\nu_k}\,
  \gamma_{h_i\ell},
  \qquad
  P^{\IO}_{\ell k}=\prod_{k'\in\cK}P_{\ell k'}^{\,\omega_{k'k}},
  \label{eq:cost}
\end{equation}
where $w_\ell$ is the wage, $\nu_k\in[0,1)$ is the share of intermediates in costs,
and $\omega_{k'k}\ge0$, with $\sum_{k'}\omega_{k'k}=1$, is the share of sector $k$'s
intermediate spending that falls on sector $k'$ \citep{CaliendoParro2015}; the
constant of the Cobb--Douglas cost function is absorbed into $\varphi_i$. The term
$\gamma_{h\ell}\ge1$, with $\gamma_{hh}=1$, is the efficiency loss from producing away
from headquarters \citep{RamondoRodriguezClare2013,HeadMayer2019}. Intermediates are
bought in the markets $(\ell,k')$ of the country where the plant operates, which
include imported varieties and are subject to that country's tariffs, so the entry of
a plant into a market lowers the input costs of every producer located in that country.

\paragraph{Head-office capability.} The productivity of plant $i$ combines a location,
a plant and a group component,
\begin{equation}
  \varphi_i=z_\ell\,e^{u_i}\times
  \begin{cases}
    \exp\bigl[(1+\delta^P\alpha_k)\,x_g\bigr] & \text{if } g\in\cG^P,\\[2pt]
    \exp\bigl[-\delta^S\alpha_k\bigr] & \text{if } g\in\cG^S,\ h(g)\in\cN^P,\\[2pt]
    1 & \text{if } g\in\cG^S,\ h(g)\notin\cN^P,
  \end{cases}
  \label{eq:productivity}
\end{equation}
where $z_\ell>0$ is the productivity of location $\ell$, common to all plants there
whatever their owner; $u_i$ is an idiosyncratic plant draw, normal with mean zero and
independent across plants; $x_g\ge0$ is the capability of parent group $g$, with
$e^{x_g}$ Pareto with unit lower bound and shape $\vartheta$; and $\alpha_k\in[0,1]$
indexes the head-office intensity, or complexity, of sector $k$
\citep{Antras2003,AntrasHelpman2004}. Head-office services (management, technology,
brands) are a non-rival capability of the group that all of its plants share, so they
enter productivity rather than the cost of inputs. A stand-alone firm has no access to
them, and its productivity falls short by a factor that grows with complexity at rate
$\delta^S\ge0$. A parent group's capability is worth more in complex sectors: its
group component is Pareto with shape $\vartheta/(1+\delta^P\alpha_k)$, with a thicker
tail where complexity is higher. Stand-alone firms in countries outside $\cN^P$ stand in for
all supply from those countries, multinational or not, and carry no penalty. The
complexity index $\alpha_k$ appears only in \eqref{eq:productivity}.\footnote{The
common alternative treats head-office services as labor bought at the parent's wage,
so that the labor term in \eqref{eq:cost} becomes
$(w_{h}^{\alpha_k}w_\ell^{1-\alpha_k})^{1-\nu_k}$ \citep{Antras2003}. This can overturn
the substitutability of labor across countries on which uniqueness of wages rests
(Appendix~\ref{app:hqinput}).}

\paragraph{Pool of potential groups.} The set of potential groups is drawn once,
before any decision is taken, and the equilibrium is computed for a given draw. Each
country--sector pair $(n,k)$ has a given number $M^S_{nk}$ of potential stand-alone
firms, and each sector $k$ has a given number $M^P_k$ of potential parent groups. The
headquarters country of each potential parent group is drawn independently from
$\cN^P$ with probability proportional to $L_hz_h^{\zeta}$, with $\zeta\ge0$, and its
capability $x_g$ is drawn from the same distribution whatever its headquarters. Larger
and more productive countries therefore host more potential parents, and their best
parents tend to be better only because more draws are taken there
\citep{GaubertItskhoki2021}.

\paragraph{Trade.} Shipping goods from $\ell$ to $n$ entails an iceberg cost
$d_{\ell n}\ge1$, with $d_{nn}=1$. Trade costs run from the country where the
good is produced, not from headquarters, so a single plant can serve several
destinations, including third markets, as an export platform \citep{Tintelnot2017}.
Country $n$ levies a tariff $t_{\ell nk}>-1$ on sector-$k$ goods produced in $\ell$,
with $t_{nnk}=0$; negative values are import subsidies. Tariffs depend on where a good
is produced, not on who owns the plant that produces it, so a country's tariffs also
fall on the foreign affiliates of the groups its residents own. We model the tariff as
a proportional wedge on the cost of delivering the good. The delivered marginal cost of
plant $i$ in market $(n,k)$ is
\begin{equation}
  c_{i,n}=a_i\,d_{\ell n}\,(1+t_{\ell nk}),\qquad \ell=\ell(i),
  \label{eq:delivered}
\end{equation}
and country $n$ collects $t_{\ell nk}\,a_i d_{\ell n}q_{i,n}$ on the sales of plant
$i$.\footnote{A duty on the exporter's price yields the same prices when all the plants a
group uses to serve $n$ face the same tariff. When they do not, it makes the group's
markup differ across its plants, and the property that a group's markup depends only on
its total share (Section~\ref{sec:market}) is lost.}

\paragraph{Multinational corporations and entry.} A plant that serves market $(n,k)$
pays a fixed cost
\begin{equation}
  F_{i,n}=f_k\,w_\ell\,f^d_{\ell n},
  \label{eq:fixed}
\end{equation}
in units of labor of the country where the plant is located, with $f_k>0$ and
$f^d_{\ell n}\ge1$ raising the cost of serving distant markets. The cost is incurred
once per market served, whatever the quantity sold, and stands for registration,
distribution and marketing; it does not depend on who owns the plant. Because it is
paid in labor, entry uses real resources, appears as wage income, and leaves the model
homogeneous of degree one in wages. We abstract from a fixed cost of setting up a
plant. Marginal costs are constant and every fixed cost is specific to a plant--market
pair, so a group's problem separates across destinations: in each market $(n,k(g))$
the group chooses the subset $T_{g,n}\subseteq\cP_g$ of its plants that serve $n$. A
parent group can serve a market from a plant at home, from a plant in the destination,
from a plant in a third country, or from several at once; a stand-alone firm decides
only whether to serve $n$. A plant operates if it serves at least one market, so which
affiliates exist, and where each of them sells, is an equilibrium outcome. This is a
simplified version of the multi-plant location problem of
\citet{Yang2023}.\footnote{A plant-level setup cost would link a group's choices across markets and
turn its problem into a combinatorial one over locations and markets jointly
\citep{Tintelnot2017}.}

\paragraph{Profit repatriation.} Country $n$ owns a share $\theta_{gn}\ge0$ of group
$g$, with $\sum_n\theta_{gn}=1$. Groups maximize their own profits whoever owns them:
ownership determines where profits are spent, not how groups price or which markets
they enter. The profit of a group net of fixed costs,
$\Pi_g=\sum_n\bigl[\pi_{g,n}-\sum_{i\in T_{g,n}}F_{i,n}\bigr]$, where $\pi_{g,n}$ is
its variable profit in market $(n,k(g))$ (Section~\ref{sec:market}), is paid to its
owners in proportion to $\theta_{gn}$, wherever it is earned. Tariff revenue is rebated
lump sum to the household of the country that collects it. Income in country $n$ is
\begin{equation}
  X_n=w_nL_n+\sum_{g\in\cG}\theta_{gn}\,\Pi_g+T_n,
  \qquad
  T_n=\sum_{k}\sum_{i\in\cA_{nk}}t_{\ell(i)nk}\,a_i\,d_{\ell(i)n}\,q_{i,n}.
  \label{eq:income}
\end{equation}
Profits are earned where groups sell and received where their owners reside, and
$\theta$ links the two. A natural benchmark assigns each group to the
country of its ultimate parent, $\theta_{gn}=\mathbf 1\{n=h(g)\}$; nothing in the
definition of equilibrium requires this form. Since a group can always choose to serve
no market, its net profit is non-negative in equilibrium. Because income includes the
profits of the groups a country owns, wherever they produce and sell, trade need not
balance country by country.

\paragraph{Timing.} The model is static. The pool of potential groups and all
productivity draws are realized first. Then, simultaneously, groups choose which of
their plants serve each market and how much each sells, and households choose their
consumption. Groups take wages, input prices and market sizes as given. They take into
account their effect on a market's price index when they set quantities, but take the
competitiveness of the market as given when they choose which plants operate
(Section~\ref{sec:market}).

\subsection{Market Structure}
\label{sec:market}

Fix a market $(n,k)$, take the set of plants that serve it as given until the last
paragraph of this section, and suppress the market's indices: we write $\sigma$, $E$
and $P$ for $\sigma_k$, $E_{nk}$ and $P_{nk}$, $c_i$ for the delivered cost $c_{i,n}$,
and $T_g$ for the plants of group $g$ that serve the market.

\paragraph{Demand.} Let $A\equiv\sum_{i\in\cA}p_i^{1-\sigma}=P^{1-\sigma}$. By
\eqref{eq:demand}, the revenue share of plant $i$ and the share of group $g$ are
\begin{equation}
  s_i=\frac{p_iq_i}{E}=\frac{p_i^{1-\sigma}}{A},
  \qquad
  S_g=\sum_{i\in T_g}s_i,
  \qquad
  \sum_g S_g=1 ,
  \label{eq:shares}
\end{equation}
where $S_g=0$ for a group that does not serve the market. The aggregate $A$ summarizes
how competitive the market is: it is large when rivals charge low prices, and then any
given price wins a small share.

\paragraph{Competition.} Groups are the strategic agents. Each group chooses the
quantities of all of its plants that serve the market, taking as given the quantities
of other groups and the overall demand for the sector, which is the level of spending
$E$ when $\eta=1$. It maximizes its variable profit $\pi_g=\sum_{i\in T_g}(p_i-c_i)q_i$.

\begin{proposition}[Markups]\label{prop:markup}
Under quantity competition, the first-order conditions of group $g$ are
$p_i=\mu(S_g)\,c_i$ for every $i\in T_g$, with
\begin{equation}
  \frac{1}{\mu(S)}=1-\frac{1-S}{\sigma}-\frac{S}{\eta}.
  \label{eq:markupC}
\end{equation}
Under price competition the markup is $\mu(S)=\varepsilon(S)/[\varepsilon(S)-1]$ with
$\varepsilon(S)=\sigma-(\sigma-\eta)S$.
\end{proposition}

\noindent Both are the markups of \citet{AtkesonBurstein2008}, evaluated at the share
of the group rather than of the plant, and both satisfy the following property whenever
$\sigma>\eta$.

\begin{assumption}\label{as:mu}
$\mu$ is continuously differentiable and strictly increasing on $[0,1)$, with
$\mu(0)=\sigma/(\sigma-1)$.
\end{assumption}

\noindent The results for a single market and for input prices use only
Assumption~\ref{as:mu}; the others use the specific form of the markup, and we say so
where they do. We take quantity competition as the baseline.\footnote{When $\eta=1$ the
markup \eqref{eq:markupC} is unbounded as $S\to1$, so a market served by a single group
has no equilibrium. Such a configuration is never an entry equilibrium if a second group can serve the
market and its spending exceeds that group's fixed cost, so we consider only
configurations with at least two active groups in every market.}

\paragraph{Cannibalization across multinational affiliates.} When a group expands one
of its plants, the sector price index falls and every variety in the market loses
sales, the group's other plants included. A single-plant firm ignores the losses it
imposes on others; a group counts the losses it imposes on itself, and these grow with
its share. A group with a negligible share competes only with rivals, at the
within-sector elasticity $\sigma$. A group holding the whole market is the sector, and
competes only with other sectors, at the elasticity $\eta$.
Equation~\eqref{eq:markupC} weights the two by the group's share, so the markup rises
from $\sigma/(\sigma-1)$ for a negligible group toward the monopoly markup. Two
implications follow. First, all plants of a group charge the same markup in a given
market, and differences in their costs show up in quantities rather than in markups.
Second, market power is a property of the group: concentration should be measured at
the level of groups, and a group charges more than its plants would charge as
independent firms. Within a group, \eqref{eq:shares} and the common markup give
\begin{equation}
  \frac{s_i}{S_g}=\frac{c_i^{1-\sigma}}{K_g},
  \qquad
  K_g\equiv\sum_{i\in T_g}c_i^{1-\sigma}.
  \label{eq:strength}
\end{equation}
We call $K_g$ the \emph{strength} of the group in the market: an index of its cost
efficiency that rises with the number of plants serving the market and with how cheaply
each delivers. Substituting $p_i=\mu(S_g)c_i$ into \eqref{eq:shares},
\begin{equation}
  \psi(S_g)=\frac{K_g}{A},
  \qquad
  \psi(S)\equiv S\,\mu(S)^{\sigma-1}.
  \label{eq:sharemap}
\end{equation}
Under Assumption~\ref{as:mu} the map $\psi$ is strictly increasing, so
$S_g=\psi^{-1}(K_g/A)$: the share of a group depends only on its strength relative to
the competitiveness of the market. However many plants a group runs, and in however
many countries, its position in the market is summarized by one number, $K_g$, and all
it needs to know about its rivals is $A$. Given the active plants, the equilibrium of the
market is the value of $A$ at which shares add up to one,
\begin{equation}
  \sum_{g}\psi^{-1}\!\left(\frac{K_g}{A}\right)=1 .
  \label{eq:clearing}
\end{equation}

\paragraph{Profits.} Variable profit is $\pi_g=E\,S_g\,m(S_g)$, where
$m(S)\equiv1-1/\mu(S)$ is the Lerner index. Under quantity competition
\eqref{eq:markupC} makes the Lerner index linear in the share,
$m(S)=(1+\kappa S)/\sigma$ with $\kappa\equiv\sigma/\eta-1$, so that
\begin{equation}
  \pi_g=\frac{E}{\sigma}\bigl(S_g+\kappa S_g^2\bigr).
  \label{eq:profit}
\end{equation}
Under price competition the Lerner index $1/\varepsilon(S)$ is not linear, and
$\kappa=(\sigma-\eta)/\sigma$ is the slope of its first-order expansion. Profits rise
more than proportionally with the share, because a larger group both sells more and
earns more on each unit sold. The curvature $\kappa$ governs the ownership results of
Section~\ref{sec:cs}. It is zero with constant markups. Under quantity competition the
Lerner index is linear, which makes those results exact; this is the main reason for our
baseline.

\paragraph{Entry.} A group's plant set matters for its profit only through its
strength. Let $K(T)=\sum_{i\in T}c_i^{1-\sigma}$ and $\Phi(T)=\sum_{i\in T}F_{i,n}$ be
the strength and the fixed cost of a set $T\subseteq\cP_g$. Given $A$, serving the
market with $T$ yields
\begin{equation}
  U_g(T;A)=\Omega\bigl(K(T);A\bigr)-\Phi(T),
  \qquad
  \Omega(K;A)\equiv E\,\psi^{-1}(K/A)\,m\bigl(\psi^{-1}(K/A)\bigr).
  \label{eq:entrypayoff}
\end{equation}
A group with $|\cP_g|$ potential plants chooses among $2^{|\cP_g|}$ sets. Because
$\Omega$ is strictly increasing in $K$, a set can be optimal only if no other set
delivers at least as much strength at no greater fixed cost. These efficient sets form
a chain, ordered by strength and by cost together, so the group's choice amounts to how
far to climb a ladder that runs from small and cheap to large and expensive. A
stand-alone firm has a ladder with a single rung. When a group's plants pay the same
fixed cost to serve the market, the efficient sets are nested: each rung adds the plant
with the next-lowest delivered cost, and the group's choice is a cutoff on its own list
of plants, the within-group analogue of the cutoff in \citet{GaubertItskhoki2021}. When
fixed costs differ across its plants, the rungs need not be nested, but the ladder
remains ordered by strength and cost.

\begin{definition}[Entry equilibrium in a market]\label{def:entry}
Given delivered costs $\{c_i\}$, fixed costs $\{F_{i,n}\}$ and spending $E$, an entry
equilibrium is a collection of active sets $\{T_g\}$ and an aggregate $A$ such that
(i) each $T_g$ maximizes $U_g(\cdot\,;A)$ over subsets of $\cP_g$, and (ii) $A$ solves
\eqref{eq:clearing} with $K_g=K(T_g)$.
\end{definition}

\noindent Groups thus evaluate entry at the prevailing degree of competition, as in a
free-entry condition, while taking into account their effect on $A$ when they set
quantities. The difference from a notion in which groups anticipate how their entry
changes $A$ matters only for groups with large shares. This notion turns each group's
entry problem into the choice of one number given one number, and this is what delivers
uniqueness in Section~\ref{sec:ge}. Appendix~\ref{app:nash} discusses the alternative.

\subsection{General Equilibrium}
\label{sec:ge}

\paragraph{Market sizes.} Let $V_{i,n'}\equiv a_i\,d_{\ell(i)n'}\,q_{i,n'}$ be the
variable cost of plant $i$ in serving country $n'$, exclusive of tariffs. By
Proposition~\ref{prop:markup}, $V_{i,n'}=r_{i,n'}/[\mu_{g(i)}(1+t_{\ell(i)n'k})]$, where
$r_{i,n'}$ is the plant's revenue in that market and $\mu_{g(i)}$ its group's markup
there. Spending on market $(n,k)$ is final spending by households plus intermediate
purchases by every plant located in $n$, wherever it sells:
\begin{equation}
  E_{nk}=\frac{\beta_kP_{nk}^{1-\eta}}{P_n^{1-\eta}}\,X_n
  +\sum_{k'\in\cK}\omega_{kk'}\,\nu_{k'}\sum_{n'\in\cN}
   \sum_{i\in\cA_{n'k'}:\,\ell(i)=n}V_{i,n'}.
  \label{eq:spending}
\end{equation}

\paragraph{Labor market.} In each country, labor demanded for production and for
serving markets equals the labor supplied:
\begin{equation}
  w_nL_n=\sum_{k\in\cK}(1-\nu_k)\sum_{n'\in\cN}\sum_{i\in\cA_{n'k}:\,\ell(i)=n}V_{i,n'}
  \;+\;\sum_{n'\in\cN}\sum_{i:\,\ell(i)=n,\ i\in T_{g(i),n'}}F_{i,n'}.
  \label{eq:labor}
\end{equation}

\begin{definition}[Equilibrium]\label{def:eq}
An equilibrium consists of wages $\{w_n\}$, price indices $\{P_{nk}\}$, market sizes
$\{E_{nk}\}$, incomes $\{X_n\}$, active sets $\{T_{g,n}\}$, and prices and quantities
$\{p_{i,n},q_{i,n}\}$ such that households maximize \eqref{eq:utility}; in every market,
prices satisfy Proposition~\ref{prop:markup}, quantities satisfy \eqref{eq:demand}, and
the active sets and the aggregate form an entry equilibrium (Definition~\ref{def:entry});
costs \eqref{eq:cost}--\eqref{eq:delivered} are evaluated at the price indices
\eqref{eq:priceindex} that they generate; market sizes and incomes satisfy
\eqref{eq:spending} and \eqref{eq:income}; and labor markets clear, \eqref{eq:labor}.
\end{definition}

\noindent All prices, costs, fixed costs and incomes are homogeneous of degree one in
wages, and shares are homogeneous of degree zero. Since variable profits and fixed
costs scale together, entry decisions are unchanged when all wages are scaled by a
common factor, so wages are determined up to scale. By Walras' law one labor-market
condition is implied by the others, and we set $w_1=1$.

\paragraph{Existence and uniqueness.} We proceed layer by layer: one market given its
active plants; entry into one market; input prices, market sizes and incomes given
wages and active plants; and wages given active plants.

\begin{proposition}[One market]\label{prop:market}
Under Assumption~\ref{as:mu}, for any active set containing at least two groups and
any positive costs, the market has exactly one equilibrium.
\end{proposition}

\noindent The proof uses \eqref{eq:clearing}: each $\psi^{-1}(K_g/A)$ falls as $A$
rises, so the sum of shares is a strictly decreasing function of a single number and
equals one exactly once. Each group's profit is strictly concave in its own quantities,
so the first-order conditions identify the unique best response.

Entry is harder. Each entry lowers the value of entering for everyone else, and
discrete games of this kind typically have several equilibria. Ranking entrants by cost
restores uniqueness only when every entrant is a separate firm, and fixing an order of
moves does so only by selecting one equilibrium. We instead separate the effect of a
group's strength on its own payoff from the effect of its rivals. Write $\Omega_K$ and
$\Omega_{KA}$ for the partial derivatives of $\Omega$.

\begin{theorem}[Entry]\label{thm:entry}
Consider quantity competition with $\eta=1$ and $\sigma\ge2$.
\begin{enumerate}[label=(\alph*)]
\item $\Omega(K;A)$ is strictly concave in $K$ at every share.
\item $\Omega_{KA}<0$ whenever the group's share is at most one half.
\item Suppose that, for any two candidate equilibria, no group would hold more than
half of the market if it operated the plants it chooses in one of them while facing the
aggregate of the other, and that no group is indifferent between two rungs of its
ladder. Then a market has at most one entry equilibrium.
\end{enumerate}
\end{theorem}

\noindent Part (a) says that strength has diminishing returns for the group, and part
(b) that it is worth less in a more competitive market. Together they imply that each
group's optimal strength falls as $A$ rises, while $A$ rises with every group's strength,
so two equilibria with different aggregates cannot both exist. No order of entry is
imposed, and groups still internalize competition among their own plants. The
hypothesis, that no group holds more than half of a market, can be checked, and it is
sufficient rather than necessary: at a share of one half the expression whose sign
gives (b) still equals $1/\sigma$. Numerical checks suggest that it also suffices under
price competition and for $\eta>1$, but we have no proof. Because entry is chosen in
whole plants, an exact equilibrium in pure strategies need not exist
(Appendix~\ref{app:nash}).

\begin{proposition}[Prices and incomes]\label{prop:prices}
Fix wages and active sets. (a) The price indices are the unique fixed point of
\eqref{eq:priceindex} combined with \eqref{eq:cost}--\eqref{eq:delivered}, and
iteration converges to it at rate $\max_k\nu_k$. (b) If tariffs are non-negative, market
sizes and incomes are the unique solution of the linear system \eqref{eq:income},
\eqref{eq:spending}.
\end{proposition}

\begin{proposition}[Existence of wages]\label{prop:exist}
Fix active sets, and suppose that $\nu_k=0$, $\eta=1$ and tariffs are non-negative.
Suppose also that (i) no set of countries is cut off from the rest: for every proper
subset $Z$ of countries, some plant located in $Z$ is active in a market of a country
outside $Z$; and (ii) fixed costs are small enough that every country's income is at
least a fraction $\epsilon>0$ of its labor income at all wages. Then there exist wages
that clear every labor market.
\end{proposition}

\noindent Uniqueness of wages is the last layer and the most delicate. Let
$W_n(w)$ be the wage bill of country $n$ implied by \eqref{eq:labor} at wages $w$, and
$M_{nj}\equiv\partial\ln W_n/\partial\ln w_j$. Homogeneity makes every row of $M$ sum
to one. Labor in different countries is a gross substitute when $M_{nj}>0$ for
$n\neq j$: a rise in any other country's wage raises country $n$'s wage bill. The
next lemma is the building block, and it holds without conditions.

\begin{lemma}[Cost pass-through]\label{lem:pass}
Fix active sets. At any wage vector, the matrix $\partial\ln P/\partial\ln w$, with
rows indexed by markets, is non-negative and its rows sum to one.
\end{lemma}

\noindent In the proof, each response is a weighted average of responses one step
down: a plant's cost averages wages and input prices, and a market's price index
averages the cost changes of the groups in it, with weights that variable markups
change but cannot make negative.

The model's solution algorithm adjusts wages by the damped tâtonnement
\[
  \Psi(w)_n=w_n^{1-\lambda}\bigl(W_n(w)/L_n\bigr)^{\lambda},\qquad\lambda\in(0,1),
\]
which raises a country's wage when its labor is in excess demand and whose fixed
points, up to scale, are the equilibrium wage vectors. Let
$d_H(w,w')=\max_n\ln(w_n/w'_n)-\min_n\ln(w_n/w'_n)$ denote Hilbert's projective
metric, which is unchanged when either vector is rescaled.

\begin{theorem}[Wages]\label{thm:wages}
Fix active sets. Let $D$ be a compact set of wage vectors that is convex in logarithms, and
suppose that (i) $M_{nj}(w)>0$ for all $n\neq j$ and every $w\in D$, and (ii)
$\lambda\,[1-M_{nn}(w)]<1$ for every $n$ and every $w\in D$. Then there is $\tau<1$ such that
$d_H(\Psi(w),\Psi(w'))\le\tau\,d_H(w,w')$ for all $w,w'\in D$. Hence $D$ contains at
most one equilibrium wage vector up to scale, and if $\Psi$ maps $D$ into itself the
algorithm converges to it from any starting point in $D$.
\end{theorem}

\noindent Since every row of $M$ sums to one, gross substitutes implies $M_{nn}<1$.
When labor demand is elastic, $M_{nn}<0$ and the undamped map ($\lambda=1$) is not
order-preserving; condition (ii), which comes from the model, says how much damping
restores the property. Uniqueness itself needs only condition (i): by a classical
argument, gross substitutes on $D$ rules out a second equilibrium in $D$
(Appendix~C.6). The next result gives conditions under which (i) holds.

\begin{proposition}[Gross substitutes]\label{prop:gs}
Suppose $\nu_k=0$ and $\eta=1$, and that no group's share exceeds one half in any market
under quantity competition. Then for $n\neq j$,
\begin{equation}
  M_{nj}\;\ge\;\bar E^{(j)}_n\equiv\sum_{i,n'}b^n_{i,n'}\,
  \frac{\partial\ln E_{n'k(i)}}{\partial\ln w_j},
  \label{eq:gsbound}
\end{equation}
where $b^n_{i,n'}$ is the share of country $n$'s wage bill paid by plant $i$ out of
its variable costs in country $n'$. Labor is a gross substitute whenever
$\bar E^{(j)}_n>0$.
\end{proposition}

\noindent The bound says that when production uses only local labor, a rise in a
foreign wage can lower country $n$'s wage bill only by shrinking the markets in which
its plants sell; every reallocation of sales between plants and between groups
favors country $n$. The one-half condition is the same as in
Theorem~\ref{thm:entry}: it signs the effect of variable markups on factor
demand. With intermediates, the same bound holds provided the cost of each plant paying
country $n$ rises by no more than the smaller of its group's average cost and the
market's price index, a condition that is vacuous when $\nu_k=0$.

These results have three limits. First, each layer is unique given the layers it
takes as given, which does not make the joint equilibrium of wages and active sets
unique: active sets change discretely with wages, and Theorem~\ref{thm:wages} holds them
fixed. Second, existence is not guaranteed everywhere: an exact entry equilibrium in
whole plants may fail to exist, and with intermediate inputs we have not proved that
market-clearing wages exist. Third,
gross substitutes is not a theorem of the model. It can fail when head-office services
are a factor input (Appendix~\ref{app:hqinput}), and with capability alone the demand
term in \eqref{eq:gsbound} can still turn negative, because a country whose groups lose
profits when a foreign wage rises spends less. That income effect is the ownership
channel the model is built to study. The hypotheses of Theorems~\ref{thm:entry} and
\ref{thm:wages} and of Proposition~\ref{prop:gs} can be checked at the computed solution, and we report these checks with the quantification.

\paragraph{Solution algorithm.} The layers above map into nested loops. Each market
is solved for $A$ from \eqref{eq:clearing}; with entry, each group's best response at a
given $A$ is found by scanning its ladder. Given wages and active sets, price indices
come from iterating the contraction of Proposition~\ref{prop:prices} and incomes from
one linear solve, and wages are updated by the tâtonnement $\Psi$ with a step that
satisfies condition (ii). An outer loop lets the groups with the largest gains change
their plant sets until no group wishes to move. This loop is not guaranteed to converge
and the joint equilibrium is not known to be unique, so we solve from several starting
points (Appendix~\ref{app:algo}).

\subsection{Comparative Statics}
\label{sec:cs}

\paragraph{Incidence within a market.} Consider a market with given active sets and
spending, and a change in delivered costs. Let $\hat c_g=\sum_{i\in T_g}(s_i/S_g)\,
d\ln c_i$ be the average cost change of group $g$, $\varepsilon_\mu(S)\equiv
S\mu'(S)/\mu(S)$ the elasticity of the markup, and
\begin{equation}
  \Lambda_g\equiv\frac{1}{1+(\sigma-1)\,\varepsilon_\mu(S_g)},
  \qquad
  \varpi_g\equiv\frac{S_g\Lambda_g}{\sum_h S_h\Lambda_h}.
  \label{eq:incidenceweights}
\end{equation}
Under quantity competition with $\eta=1$, $\Lambda_g=(1-S_g)/[1+(\sigma-2)S_g]$, which
falls from one toward zero as the share rises.

\begin{proposition}[Incidence]\label{prop:incidence}
Holding active sets and spending fixed,
\begin{equation}
  d\ln S_g=-(\sigma-1)\Lambda_g\bigl(\hat c_g-d\ln P\bigr),
  \qquad
  d\ln P=\sum_g\varpi_g\,\hat c_g .
  \label{eq:incidence}
\end{equation}
If only group $j$'s costs change, its sales-weighted average price rises by the fraction
$\rho_j=1-(\sigma-1)\varepsilon_\mu(S_j)\Lambda_j(1-\varpi_j)\in(0,1]$ of its average
cost increase $\hat c_j$, and every rival $h$ gains share,
$d\ln S_h=(\sigma-1)\Lambda_h\varpi_j\hat c_j$, and raises its price.
\end{proposition}

\noindent In the limit of small groups, which is monopolistic competition,
$\Lambda_g=1$, $\varpi_g=S_g$ and $\rho_j=1$. Large groups instead absorb part of a cost
increase in their markup, so pass-through falls with the group's share, and since
$\varpi_g<S_g$ for the larger groups, the price index moves by less than their share
would suggest. Rivals are not passive either: a cost increase for one group raises their
shares, markups and prices. Because tariffs depend on where goods are produced, a tariff on goods from
$\ell$ raises $\hat c_g$ only by the share of group $g$'s sales that comes from its
plants in $\ell$: a group able to serve the market from several countries is partly
shielded, while a stand-alone exporter from $\ell$ bears the tariff in full.

\paragraph{Ownership and the level of profit income.} Consider a country $H$ with
ownership shares $\theta_g\equiv\theta_{gH}$, and its profit income from a market,
$\pi^H=\sum_g\theta_g\pi_g$.

\begin{proposition}[Ownership-weighted concentration]\label{prop:ownership}
Under quantity competition,
\begin{equation}
  \frac{\pi^H}{E}=\frac{1}{\sigma}\Bigl[\bar\theta
  +\kappa\bigl(\bar\theta\cdot\mathrm{HHI}+\mathrm{Cov}_S(\theta,S)\bigr)\Bigr],
  \label{eq:decomp}
\end{equation}
where $\bar\theta=\sum_gS_g\theta_g$ is the sales-weighted ownership share,
$\mathrm{HHI}=\sum_gS_g^2$, and
$\mathrm{Cov}_S(\theta,S)=\sum_g\theta_gS_g^2-\bar\theta\cdot\mathrm{HHI}$ is the
covariance between ownership and share computed with sales weights. A calculation that
applies $\bar\theta$ to total profit in the market, $\bar\theta\sum_g\pi_g$, misses
exactly $(E\kappa/\sigma)\,\mathrm{Cov}_S(\theta,S)$.
\end{proposition}

\noindent The first term is what a model with constant markups delivers: ownership
times revenue over $\sigma$. The second is an ownership-weighted Herfindahl index. Its
component $\bar\theta\cdot\mathrm{HHI}$ can be built from a country's ownership share and
a concentration index; the covariance cannot, since it requires knowing which groups a
country owns and how large each one is. Because profits are convex in shares, two
countries with the same share of a market's sales earn different profits if one owns
the large groups and the other the small ones. With constant markups, as under
monopolistic competition, $\kappa=0$ and a country's profit income is its ownership
share times total profit, whatever the distribution of ownership across groups; that
benchmark lies outside the oligopoly model, where $\kappa=\sigma/\eta-1>0$. The
decomposition is exact under quantity competition and a first-order approximation under
price competition. It concerns variable profit, before the fixed costs that a country's
groups pay to serve the market. Summed over markets with spending weights, the same
algebra splits the economy-wide correction into a covariance across markets between
ownership shares and concentration, and the average of the covariances within markets.

\paragraph{Ownership and the response of profit income.} Welfare depends on how profit
income responds to policy, and the response is governed by a different statistic.

\begin{proposition}[Response of owned profits]\label{prop:response}
Under quantity competition, holding active sets and spending fixed, a
change in the delivered costs of group $j$ alone changes country $H$'s profit income by
\begin{equation}
  \frac{d\pi^H}{E}=\frac{\sigma-1}{\sigma}\,\varpi_j\,\hat c_j
  \sum_{h\neq j}S_h\Lambda_h\Bigl[\theta_h\bigl(1+2\kappa S_h\bigr)
  -\theta_j\bigl(1+2\kappa S_j\bigr)\Bigr].
  \label{eq:response}
\end{equation}
\end{proposition}

\noindent A cost increase for group $j$ alone, for example a tariff on goods from a
country from which only $j$ serves the market, reallocates sales from $j$ to its rivals
in proportion to $S_h\Lambda_h$, and each unit of share is worth $\theta(1+2\kappa S)$ to
country $H$, up to the factor $E/\sigma$. The country gains when the rivals it owns are
worth more to it at the margin than the group whose costs rose. The weights are
$S_h\Lambda_h$ rather than $S_h^2$, so large groups count for less in the response than
in the level, and what matters is the ownership of the affected group relative to its
rivals, not the covariance that governs the level. The ownership-weighted Herfindahl
index measures a country's stake in a market but is not a sufficient statistic for how
that stake responds to a tariff; in general equilibrium the response also runs through
wages, entry and market sizes.

\paragraph{Entry.} The proof of Theorem~\ref{thm:entry} also gives comparative statics.

\begin{corollary}[Entry responses]\label{cor:entry}
Under quantity competition with $\eta=1$, the strength a group chooses is
non-decreasing in spending $E$, holding $A$ fixed. It is non-increasing in $A$, holding
$E$ fixed, provided the group's share would not exceed one half with any of its rungs at
the aggregates compared.
\end{corollary}

\noindent A more competitive market leads a group to climb fewer rungs of its ladder,
and a larger market to climb more. Since every rung is a set of plants, these responses
determine which affiliates operate, and through \eqref{eq:spending} and \eqref{eq:labor}
they feed back into market sizes and wages.

\paragraph{Wages and prices.} By Lemma~\ref{lem:pass}, input--output links and markups
propagate a wage change across price indices but cannot reverse or amplify it. Who gains from a fall in $j$'s wage depends on ownership. The fall lowers
the costs of every group with plants in $j$, and part of the gain accrues as profit to
the owners of those groups, wherever they reside. A policy that lowers a foreign wage
can therefore raise the income of the country that owns the plants located there. This
is the general-equilibrium channel through which ownership matters for trade policy.

\setlength{\bibsep}{1pt}
\begin{thebibliography}{99}

\bibitem[Antr\`as(2003)]{Antras2003}
Antr\`as, Pol. ``Firms, Contracts, and Trade Structure.'' \emph{Quarterly Journal of
Economics} 118(4): 1375--1418, 2003.

\bibitem[Antr\`as and Helpman(2004)]{AntrasHelpman2004}
Antr\`as, Pol, and Elhanan Helpman. ``Global Sourcing.'' \emph{Journal of Political
Economy} 112(3): 552--580, 2004.

\bibitem[Arrow, Block, and Hurwicz(1959)]{ArrowBlockHurwicz1959}
Arrow, Kenneth J., H. D. Block, and Leonid Hurwicz. ``On the Stability of the
Competitive Equilibrium, II.'' \emph{Econometrica} 27(1): 82--109, 1959.

\bibitem[Atkeson and Burstein(2008)]{AtkesonBurstein2008}
Atkeson, Andrew, and Ariel Burstein. ``Pricing-to-Market, Trade Costs, and
International Relative Prices.'' \emph{American Economic Review} 98(5): 1998--2031,
2008.

\bibitem[Birkhoff(1957)]{Birkhoff1957}
Birkhoff, Garrett. ``Extensions of Jentzsch's Theorem.'' \emph{Transactions of the
American Mathematical Society} 85(1): 219--227, 1957.

\bibitem[Caliendo and Parro(2015)]{CaliendoParro2015}
Caliendo, Lorenzo, and Fernando Parro. ``Estimates of the Trade and Welfare Effects of
NAFTA.'' \emph{Review of Economic Studies} 82(1): 1--44, 2015.

\bibitem[Gaubert and Itskhoki(2021)]{GaubertItskhoki2021}
Gaubert, Cecile, and Oleg Itskhoki. ``Granular Comparative Advantage.''
\emph{Journal of Political Economy} 129(3): 871--939, 2021.

\bibitem[Head and Mayer(2019)]{HeadMayer2019}
Head, Keith, and Thierry Mayer. ``Brands in Motion: How Frictions Shape Multinational
Production.'' \emph{American Economic Review} 109(9): 3073--3124, 2019.

\bibitem[Mas-Colell, Whinston, and Green(1995)]{MWG1995}
Mas-Colell, Andreu, Michael D. Whinston, and Jerry R. Green. \emph{Microeconomic
Theory}. Oxford University Press, 1995.

\bibitem[Milgrom and Shannon(1994)]{MilgromShannon1994}
Milgrom, Paul, and Chris Shannon. ``Monotone Comparative Statics.''
\emph{Econometrica} 62(1): 157--180, 1994.

\bibitem[Ramondo and Rodr\'iguez-Clare(2013)]{RamondoRodriguezClare2013}
Ramondo, Natalia, and Andr\'es Rodr\'iguez-Clare. ``Trade, Multinational Production,
and the Gains from Openness.'' \emph{Journal of Political Economy} 121(2): 273--322,
2013.

\bibitem[Seneta(2006)]{Seneta2006}
Seneta, Eugene. \emph{Non-negative Matrices and Markov Chains}. Revised printing.
Springer, 2006.

\bibitem[Tintelnot(2017)]{Tintelnot2017}
Tintelnot, Felix. ``Global Production with Export Platforms.'' \emph{Quarterly Journal
of Economics} 132(1): 157--209, 2017.

\bibitem[Yang(2023)]{Yang2023}
Yang, Chenying. ``Location Choices of Multi-plant Oligopolists: Theory and Evidence
from the Cement Industry.'' Working paper, 2023.

\end{thebibliography}

\newpage
\section*{Appendix}
\renewcommand{\theHsubsection}{app.\Alph{subsection}}
\renewcommand{\theHequation}{app.\arabic{equation}}
\setcounter{subsection}{0}
\setcounter{equation}{0}
\renewcommand{\theequation}{A.\arabic{equation}}

\renewcommand{\thesubsection}{\Alph{subsection}}

\subsection{The Market Game}
\label{app:market}

\paragraph{A.1 Markups.} Fix a market. With $\eta$ the upper-tier elasticity, write
spending as $E=D\,P^{1-\eta}$, with the shifter $D$ taken as given by firms; $\eta=1$
gives $E=D$. Let $\rho=(\sigma-1)/\sigma$ and $Q=(\sum_iq_i^{\rho})^{1/\rho}$. Inverse
demand is $p_i=D^{1/\eta}Q^{1/\sigma-1/\eta}q_i^{-1/\sigma}$, so group $g$'s revenue is
$R_g=D^{1/\eta}Q^{1/\sigma-1/\eta}\sum_{j\in T_g}q_j^{\rho}$. Using
$\partial\ln Q/\partial q_i=q_i^{\rho-1}/Q^{\rho}$ and $s_i=q_i^\rho/Q^\rho$,
\[
  \frac{\partial R_g}{\partial q_i}
  =p_i\Bigl[\rho+\Bigl(\frac1\sigma-\frac1\eta\Bigr)S_g\Bigr]
  =p_i\Bigl[1-\frac{1-S_g}{\sigma}-\frac{S_g}{\eta}\Bigr].
\]
Setting marginal revenue equal to $c_i$ gives \eqref{eq:markupC}; the bracket is common
to all $i\in T_g$ and depends only on $S_g$. Under price competition, demand is
$q_i=p_i^{-\sigma}P^{\sigma-\eta}D$, so
$\partial\ln q_j/\partial\ln p_i=-\sigma\mathbf 1\{i=j\}+(\sigma-\eta)s_i$. The
first-order condition for $p_i$ is
$q_i-\sigma(p_i-c_i)q_i/p_i+(\sigma-\eta)(s_i/p_i)\sum_{j\in T_g}(p_j-c_j)q_j=0$. Substituting
$p_j=\mu c_j$ for all $j\in T_g$, so that $(p_j-c_j)q_j=(1-1/\mu)r_j$ and
$\sum_{j\in T_g}r_j=S_gE$, it reduces to
$(1-1/\mu)\bigl[\sigma-(\sigma-\eta)S_g\bigr]=1$, which is the stated markup.

\paragraph{A.2 Proof of Proposition~\ref{prop:market}.} \emph{The share map.}
$\psi(S)=S\mu(S)^{\sigma-1}$ is the product of $S$ and of $\mu(S)^{\sigma-1}$, both
positive and strictly increasing on $(0,1)$ under Assumption~\ref{as:mu}, and
$\psi(0)=0$; so $\psi$ has a continuous, strictly increasing inverse.
\emph{Best responses.} Given rivals' quantities, group $g$'s profit is
$D^{1/\eta}X_g(X_g+B_g)^{e}-\sum_{i\in T_g}c_iq_i$, where $X_g=\sum_{i\in T_g}q_i^\rho$,
$B_g$ is the same index for rivals, and $e=(\eta-\sigma)/[\eta(\sigma-1)]\in[-1,0)$. The
function $X\mapsto X(X+B)^e$ is strictly increasing and strictly concave for $B>0$: its
derivatives are $(X+B)^{e-1}[(1+e)X+B]>0$ and $e(X+B)^{e-2}[(1+e)X+2B]<0$. The index
$X_g$ is strictly concave in $(q_i)_{i\in T_g}$. A strictly increasing, strictly
concave function of a strictly concave function is strictly concave, and cost is
linear, so profit is strictly concave in the group's own quantities. Since
$\partial X_g/\partial q_i\to\infty$ as $q_i\to0$, the optimum is interior and is the
unique solution of the first-order conditions. Under price competition, dividing the first-order condition of
A.1 by $q_i$ gives $1-c_i/p_i=[1+(\sigma-\eta)\pi_g/E]/\sigma$, whose right side does
not depend on $i$, so every interior critical point has a common markup $\mu$. With a
common markup, $S_g(\mu)=\mu^{1-\sigma}K_g/(\mu^{1-\sigma}K_g+B'_g)$, where $B'_g$ is the
sum of $p^{1-\sigma}$ over rivals, and $E=DP^{1-\eta}$, so that
\[
  \frac{d\ln\pi_g}{d\ln\mu}=\frac{1}{\mu-1}-\bigl[\varepsilon(S_g(\mu))-1\bigr].
\]
The first term falls from infinity to zero as $\mu$ rises, and the second rises because
$S_g$ falls with $\mu$; the derivative changes sign once, from positive to negative, and
the best response is unique. Pricing a plant out of the market is never optimal, since
adding it at the common markup raises the group's share without lowering its margin. \emph{The aggregate.} By
Proposition~\ref{prop:markup} and \eqref{eq:sharemap}, an equilibrium is an $A$ solving
\eqref{eq:clearing}, and conversely any such $A$ generates prices
$p_i=\mu(S_g)c_i$ that satisfy every group's first-order conditions and are consistent
with $A$, since $\sum_ip_i^{1-\sigma}=\sum_g\mu(S_g)^{1-\sigma}K_g=A\sum_gS_g=A$. The
left side of \eqref{eq:clearing}, $\mathcal S(A)$, is continuous and strictly decreasing.
As $A\to\infty$ every share tends to zero. Let $\bar\psi=\lim_{S\to1}\psi(S)$, infinite
when $\eta=1$; as $A$ falls toward $\max_gK_g/\bar\psi$ the largest group's share tends
to one while every other share stays positive, so $\mathcal S$ exceeds one. Hence
$\mathcal S(A)=1$ has exactly one solution, and every other variable is a function of it.
\qed

\subsection{Entry}
\label{app:entry}

\paragraph{B.1 Two elasticities.} Let $G(S)\equiv\pi'(S)/\psi'(S)$, where
$\pi(S)=ES\,m(S)$, and define $\varepsilon_G=SG'/G$ and $\varepsilon_\psi=S\psi'/\psi$.
Differentiating $\psi(S)=K/A$ gives $\partial S/\partial K=1/(A\psi')$ and
$\partial S/\partial A=-\psi/(A\psi')$. Hence $\Omega_K=G(S)/A>0$,
$\Omega_{KK}=G'(S)/(A^2\psi'(S))$, and
\[
  \Omega_{KA}=-\frac{1}{A^2}\Bigl[G'(S)\frac{\psi(S)}{\psi'(S)}+G(S)\Bigr]
  =-\frac{G(S)}{A^2\varepsilon_\psi(S)}\bigl[\varepsilon_G(S)+\varepsilon_\psi(S)\bigr].
\]
So $\Omega$ is strictly concave in $K$ if and only if $\varepsilon_G<0$, and
$\Omega_{KA}<0$ if and only if $\varepsilon_G+\varepsilon_\psi>0$.

\paragraph{B.2 Proof of Theorem~\ref{thm:entry}(a)--(b).} With $\eta=1$,
$\mu(S)=\sigma/[(\sigma-1)(1-S)]$. Set $E=1$, $a\equiv2(\sigma-1)$ and
$b\equiv\sigma-2\ge0$. Then $\pi'(S)=(1+aS)/\sigma$ and
$\psi'(S)=C(1-S)^{-\sigma}(1+bS)$ with $C=[\sigma/(\sigma-1)]^{\sigma-1}$, so
\[
  \varepsilon_G(S)=S\,f(S),\qquad
  f(S)=\frac{a}{1+aS}-\frac{b}{1+bS}-\frac{\sigma}{1-S}.
\]
At $S=0$, $f(0)=a-b-\sigma=0$. Since $a>b\ge0$, $a/(1+aS)>b/(1+bS)\ge0$, so
$f'(S)=-a^2/(1+aS)^2+b^2/(1+bS)^2-\sigma/(1-S)^2<0$. Hence $f<0$ and $\varepsilon_G<0$ on
$(0,1)$, which is (a). Next, $\varepsilon_\psi=(1+bS)/(1-S)$, and using
$1+bS-\sigma S=1-2S$,
\[
  \varepsilon_G+\varepsilon_\psi=\frac{aS}{1+aS}-\frac{bS}{1+bS}+\frac{1-2S}{1-S}.
\]
The map $x\mapsto x/(1+x)$ is increasing and $a>b$, so the first two terms sum to a
positive number; the third is non-negative when $S\le\frac12$. This is (b). At
$S=\frac12$ the expression equals $(\sigma-1)/\sigma-(\sigma-2)/\sigma=1/\sigma$. \qed

\paragraph{B.3 Proof of Theorem~\ref{thm:entry}(c).} \emph{Step 1: the ladder.} Since
$\Omega_K>0$, a set $T$ is optimal only if no $T'$ has $K(T')\ge K(T)$ and
$\Phi(T')\le\Phi(T)$ with one inequality strict. The remaining efficient sets can be
ordered so that $K$ and $\Phi$ rise together. \emph{Step 2: single crossing.} Take
$A<A'$ and two rungs with $K_1<K_2$, and let
$\Delta(x)=\int_{K_1}^{K_2}\Omega_K(K;x)\,dK-[\Phi_2-\Phi_1]$ be the net gain from the
larger rung at aggregate $x$. The share $\psi^{-1}(K/x)$ rises with $K$ and falls with
$x$, so on $[K_1,K_2]\times[A,A']$ it is largest at $(K_2,A)$. If
$\psi^{-1}(K_2/A)\le\frac12$, part (b) gives $\Omega_{KA}<0$ throughout, and $\Delta$
is strictly decreasing on $[A,A']$. \emph{Step 3: the aggregate.} Given strengths, the
aggregate $A(\{K_g\})$ solving \eqref{eq:clearing} is strictly increasing in each $K_g$:
raising $K_j$ at a given $A$ raises the left side, which is strictly decreasing in $A$.
\emph{Step 4.} Suppose $(A,\{K_g\})$ and $(A',\{K'_g\})$ are equilibria with $A<A'$,
and suppose $K'_g>K_g$ for some $g$. Apply Step 2 with $K_1=K_g$ and $K_2=K'_g$: by the
hypothesis, group $g$ would hold at most half of the market with the plants it chooses
at $A'$ while facing $A$, so $\psi^{-1}(K'_g/A)\le\frac12$ and $\Delta$ is strictly
decreasing on $[A,A']$. Optimality at $A'$ gives $\Delta(A')\ge0$, hence $\Delta(A)>0$,
so at $A$ the group strictly prefers $K'_g$ to $K_g$, contradicting optimality at $A$.
Hence $K'_g\le K_g$ for every $g$ \citep[the single-crossing argument of][]{MilgromShannon1994},
and by Step 3, $A'=A(\{K'_g\})\le A(\{K_g\})=A$, a contradiction. Equilibria with the
same aggregate have the same best responses, and with no ties the same active sets.
\qed

\paragraph{B.4 The notion of entry equilibrium and existence.}
\label{app:nash}
Definition~\ref{def:entry} evaluates every alternative set at the prevailing aggregate
$A$. In a two-stage game in which the market is re-solved after a deviation, a group
that adds a plant anticipates that $A$ rises, which lowers the value of the deviation;
the two notions coincide for groups with negligible shares and can differ for large
ones. Theorem~\ref{thm:entry} applies to the notion of Definition~\ref{def:entry}. For
the two-stage notion we have only numerical evidence, from markets small enough to
enumerate every configuration and re-solve the market after every deviation. Existence
in pure strategies is not guaranteed under either notion: a group's choice is discrete,
so $K^*_g(A)$ jumps, and \eqref{eq:clearing} may have no solution with each group on
a rung. The algorithm (Appendix~\ref{app:algo}) then returns the configuration with the
fewest groups wishing to deviate.

\subsection{General Equilibrium}
\label{app:ge}

\paragraph{C.1 Proof of Proposition~\ref{prop:prices}.} \emph{(a)} Let $T$ map log price
indices $\ln P$ into the log price indices generated by the costs
\eqref{eq:cost}--\eqref{eq:delivered} that $P$ implies, holding wages and active sets
fixed. Multiplying every cost in a market by a constant leaves shares and markups
unchanged and multiplies every price by that constant, so each market's price index is
homogeneous of degree one in the costs entering it; by the market step of the proof
of Lemma~\ref{lem:pass} (C.3) it is also non-decreasing in each of them. In \eqref{eq:cost}, $\ln a_i$ depends on
$\ln P$ through $\nu_k\sum_{k'}\omega_{k'k}\ln P_{\ell k'}$, with non-negative weights
summing to $\nu_k$. Hence $T$ is monotone and
$T(\ln P+x\mathbf 1)\le T(\ln P)+(\max_k\nu_k)\,x\mathbf 1$ for $x\ge0$, so by
Blackwell's sufficient conditions $T$ is a contraction in the sup norm with modulus
$\max_k\nu_k<1$. \emph{(b)} Given shares and markups, the second term of
\eqref{eq:spending} is linear in market sizes, $E=\mathcal EX+BE$, where $B\ge0$ and
the column of $B$ for market $(n',k')$ sums to $\nu_{k'}\sum_i s_i/[\mu(1+t)]$. This is
below $\nu_{k'}<1$ when tariffs are non-negative, since markups exceed one. So
$\rho(B)<1$ and $E=(I-B)^{-1}\mathcal EX\ge0$. Variable profits and tariff revenue are
linear in $E$ and, with non-negative tariffs, non-negative. Let $\Phi_g$ be group $g$'s
total fixed costs. Then \eqref{eq:income} reads $X=wL-\theta'\Phi+\mathcal MX$ with
$\mathcal M\ge0$. Summing over countries, $\mathbf 1'\mathcal MX$ equals world
spending minus intermediate purchases minus the variable wage bill, that is world final
spending minus the variable wage bill; since every unit of spending generates a positive
variable wage bill when $\nu_k<1$, each column of $\mathcal M$ sums to less than one,
$\rho(\mathcal M)<1$, and $X=(I-\mathcal M)^{-1}(wL-\theta'\Phi)$ is the unique
solution. With import subsidies the system is still linear, and it has a unique
solution whenever $I-B$ and $I-\mathcal M$ are non-singular, which can be checked
directly. \qed

\paragraph{C.2 Proof of Proposition~\ref{prop:exist}.} Normalize wages to the unit
simplex and let $z_n(w)=W_n(w)/w_n-L_n$ be excess demand for labor. By
Propositions~\ref{prop:market}--\ref{prop:prices}, $z$ is continuous on the interior of
the simplex and homogeneous of degree zero. It satisfies Walras' law
$\sum_nw_nz_n(w)=0$, which follows from adding \eqref{eq:income} across countries and
using \eqref{eq:spending} and the definition of $W_n$ as the right side of
\eqref{eq:labor}. It is bounded below by $-L$. It remains to check the boundary
condition of \citet[Proposition 17.C.1]{MWG1995}: if $w^m\to\bar w$, where
$Z=\{n:\bar w_n=0\}$ is non-empty, then $\max_nz_n(w^m)\to\infty$.

With $\nu_k=0$, a plant's delivered cost is proportional to the wage where it is
located. So the costs of plants located in $Z$ go to zero, while all other costs stay in
a compact subset of $(0,\infty)$. By (i), some plant located in $Z$ is active in a
market $(n',k)$ with $n'\notin Z$. By (ii) and $\eta=1$,
$E_{n'k}=\beta_kX_{n'}\ge\beta_k\epsilon w_{n'}L_{n'}$, which stays bounded away from
zero. For $m$ large, the lowest-cost plant $i$ in that market is located in some
$\ell\in Z$. Let $g$ be its group and $J$ the number of plants in the market. Its
variable cost bill is $V_i=E_{n'k}s_i/[\mu_g(1+t)]$, all of it paid to labor in $\ell$,
so it suffices to show that $V_i/w_\ell\to\infty$.

If $S_g\le\frac12$, then $\mu_g\le\bar\mu\equiv\mu(\frac12)$. Every price is at least
$c_i$, since markups exceed one, so $P^{1-\sigma}\le Jc_i^{1-\sigma}$ and
$s_i=(\mu_gc_i/P)^{1-\sigma}\ge\bar\mu^{1-\sigma}/J$. Then $V_i$ is bounded away from
zero while $w_\ell\to0$.

If $S_g>\frac12$, every rival $h$ has $S_h<\frac12$, so
$\psi(S_h)\le\bar\mu^{\sigma-1}S_h\le\bar\mu^{\sigma-1}(1-S_g)$, while
$\psi(S_g)\ge\frac12\mu_g^{\sigma-1}$. Under both conducts with $\eta=1$,
$\mu(S)(1-S)\le\sigma/(\sigma-1)$. Together with \eqref{eq:sharemap}, this gives
$K_h/K_g=\psi(S_h)/\psi(S_g)\le C\mu_g^{-\sigma}$ for a constant $C$. Rival strength
$K_h$ is bounded below, because all costs are bounded above. Plant $i$ is the cheapest
in its group, so $K_g\le Jc_i^{1-\sigma}$ and $s_i\ge S_g/J>1/(2J)$. Hence $V_i\ge
c\,c_i^{(\sigma-1)/\sigma}$ for some $c>0$, and since $c_i$ is proportional to $w_\ell$,
$V_i/w_\ell$ grows like $w_\ell^{-1/\sigma}$.

In either case $z_\ell\to\infty$, and the result follows from
\citet[Proposition 17.C.1]{MWG1995}. \qed

\paragraph{C.3 Proof of Lemma~\ref{lem:pass}.} Hats denote log changes in response to
a change in wages. \emph{Costs.} Trade costs and tariffs are fixed, so by \eqref{eq:cost} and
\eqref{eq:delivered}, $\hat c_i=\hat a_i=(1-\nu_k)\hat w_\ell+\nu_k\sum_{k'}\omega_{k'k}\hat P_{\ell k'}$, with non-negative coefficients
summing to one. \emph{Markets.} Differentiating $S_g=\mu(S_g)^{1-\sigma}K_g/A$, with
$\hat K_g=(1-\sigma)\hat c_g$ by \eqref{eq:strength} and $\hat A=(1-\sigma)\hat P$,
gives $\hat S_g=-(\sigma-1)\Lambda_g(\hat c_g-\hat P)$, with $\Lambda_g$, $\varpi_g$
and $\hat c_g$ as defined in \eqref{eq:incidenceweights}. Shares sum to one, so
$\sum_gS_g\hat S_g=0$, which gives $\hat P=\sum_g\varpi_g\hat c_g$ with weights
$\varpi_g>0$ summing to one. \emph{Stacking.} Substituting costs into markets gives
$\hat P=\Gamma\hat w+\mathcal N\hat P$ with $\Gamma,\mathcal N\ge0$,
$(\Gamma+\mathcal N)\mathbf 1=\mathbf 1$, and the rows of $\mathcal N$ summing to
$\nu_k<1$. Hence $\rho(\mathcal N)<1$, $(I-\mathcal N)^{-1}=\sum_{s\ge0}\mathcal
N^s\ge0$, and $\partial\ln P/\partial\ln w=(I-\mathcal N)^{-1}\Gamma\ge0$ with row sums
$(I-\mathcal N)^{-1}(\mathbf 1-\nu)=\mathbf 1$. \qed

\paragraph{C.4 A decomposition of factor payments.} Let $V_i$ be plant $i$'s variable
cost in a market $(n',k)$, with tariffs fixed. Since $V_i=r_i/[\mu_g(1+t)]$,
$\hat V_i=\hat E+\hat s_i-\hat\mu_g$, with $\hat s_i=(1-\sigma)(\hat c_i+\hat\mu_g-\hat
P)$ and $\hat\mu_g=\varepsilon_\mu\hat S_g$. Using C.3 and
$(\sigma-1)=(\sigma-1)\Lambda_g[1+(\sigma-1)\varepsilon_\mu]$,
\begin{equation}
  \hat V_i=\hat E-(\sigma-1)(\hat c_i-\hat c_g)
  -(1-\varepsilon_\mu)(\sigma-1)\Lambda_g(\hat c_g-\hat P).
  \tag{C.4}\label{eq:vdecomp}
\end{equation}
The first term is the growth of the market, the second the reallocation of sales within
the group, the third the reallocation between the group and its rivals. The coefficient
of the third term
is non-negative if and only if $\varepsilon_\mu\le1$, which under quantity competition
with $\eta=1$, where $\varepsilon_\mu=S/(1-S)$, is exactly $S_g\le\frac12$.

\paragraph{C.5 Proof of Proposition~\ref{prop:gs}.} Raise $w_j$ alone, $j\neq n$. With
$\nu_k=0$ and head-office services a capability, the whole variable labor bill of a
plant is paid in its location, so the plants paying country $n$ are located in $n$, and
their costs do not respond: $\hat c_i=0$. Fixed costs \eqref{eq:fixed} are paid at
$w_n$ and do not respond either. By Lemma~\ref{lem:pass}, $\hat c_g,\hat P\in[0,1]$. Let
$\chi=(1-\varepsilon_\mu)(\sigma-1)\Lambda_g$. The one-half condition and
$\Lambda_g\le1$ give $0\le\chi\le\sigma-1$, so the last two terms of
\eqref{eq:vdecomp} equal $(\sigma-1)\hat c_g-\chi(\hat c_g-\hat P)$. If $\hat c_g\ge\hat
P$ this is at least $(\sigma-1)\hat P$; otherwise both parts are non-negative. In either
case it is at least $(\sigma-1)\min\{\hat c_g,\hat P\}\ge0$, so $\hat V_i\ge\hat E$.
Weighting by $b^n_{i,n'}$ and adding the fixed-cost part, whose response is zero,
gives \eqref{eq:gsbound}. \qed

\paragraph{C.6 Gross substitutes and uniqueness.} The following argument is classical
\citep{ArrowBlockHurwicz1959}; see also \citet[Section 17.F]{MWG1995}. It shows that condition (i) of Theorem~\ref{thm:wages} alone gives uniqueness. Suppose
$M_{nj}>0$ for $n\neq j$ at every point of a set $D$ that is convex in logarithms.
Since $M$ is homogeneous of degree zero in wages, the same holds at every rescaling of a
point of $D$, and the set of these rescalings is still convex in logarithms. Let
$w,w^*\in D$ be non-proportional equilibria. Scale $w^*$ so that $w^*\le w$ with
equality at some $n_0$, and move from $w^*$ to $w$ along the segment in logarithms.
Every wage other than $w_{n_0}$ weakly rises and at least one strictly rises, so
$W_{n_0}$ strictly rises while $w_{n_0}$ does not. Then
$z_{n_0}(w)>z_{n_0}(w^*)=0$, a contradiction.

\paragraph{C.7 Proof of Theorem~\ref{thm:wages}.} In logarithms, $\Psi$ becomes a map
with Jacobian $(1-\lambda)I+\lambda M$, whose rows sum to one. Its off-diagonal entries
are positive by (i) and its diagonal entries by (ii), so it is a strictly positive
stochastic matrix at every point of $D$. In logarithms, $d_H$ is the oscillation
seminorm $\mathrm{osc}(x)=\max_nx_n-\min_nx_n$. By the mean value theorem in integral
form, the difference between the images of two points is $\bar A$ times the difference
between the points, where $\bar A$ is the average of the Jacobian along the segment
joining them, which stays in $D$. The contraction coefficient of a stochastic matrix in
the oscillation seminorm, $\delta(A)=\frac12\max_{i,j}\|A_{i\cdot}-A_{j\cdot}\|_1$, is
convex in $A$, so $\delta(\bar A)\le\sup_D\delta$. For a strictly positive stochastic
matrix, $\delta(A)\le\tanh(\Delta(A)/4)<1$, where
$\Delta(A)=\max_{i,j,k,l}\ln(A_{ik}A_{jl}/A_{jk}A_{il})$
\citep{Birkhoff1957,Seneta2006}. Since $D$ is compact and $M$ is continuous, $\Delta$ is
bounded on $D$, and $\tau=\sup_D\tanh(\Delta/4)<1$ gives the contraction. A fixed point
of $\Psi$ up to scale, $\Psi(w)=cw$, has $W_n=c^{1/\lambda}w_nL_n$ for every $n$. Walras'
law then forces $c=1$, so it is an equilibrium. Two such points in $D$ would contradict
the contraction. If $\Psi$ maps $D$ into itself, the iterates converge by the
contraction mapping theorem, since $D$ is compact. \qed

\paragraph{C.8 Head-office services as a factor input.}
\label{app:hqinput}
Suppose instead that the labor term in \eqref{eq:cost} is
$(w_h^{\alpha_k}w_\ell^{1-\alpha_k})^{1-\nu_k}$, so that country $h$ is paid the share
$\alpha_k$ of the labor bill of every plant its groups run. All plants of a group
report to the same headquarters, so country $h$ receives the share $\alpha_k$ of the
group's whole bill in the market. When the group's plants serving a market face the
same tariff, as in the example below, the within-group term of \eqref{eq:vdecomp}
averages to zero over the group. Consider a country $n$ whose labor income consists
only of head-office payments from groups whose plants are all located in a single
country $j$, with $\nu_k=0$. Raising $w_j$ raises those groups' costs by $\hat
c_g=1-\alpha_k$, while the price indices of the markets they serve rise only in
proportion to the groups' incidence weights. By \eqref{eq:vdecomp}, $M_{nj}$ is then a
weighted sum over these groups of $\hat E-(1-\varepsilon_\mu)(\sigma-1)\Lambda_g(\hat
c_g-\hat P)$. When the groups are small, $\Lambda_g\approx1$, $\varepsilon_\mu\approx0$
and $\hat c_g-\hat P\approx1-\alpha_k$, so $M_{nj}$ is negative whenever
$(\sigma-1)(1-\alpha_k)$ exceeds the growth of the markets they serve, although no
group is near one half of any market. The failure is worst when $\alpha_k$ is
small: a large head-office share ties the group's cost to its own country's wage and
insulates it from $w_j$. Gross substitutes fails because country $n$'s income is
concentrated on firms whose costs are set by a single foreign country. The argument
holds at any wage vector at which these groups are small, including the equilibrium
one.

\subsection{Comparative Statics}
\label{app:cs}

\paragraph{D.1 Proof of Proposition~\ref{prop:incidence}.} The first two equations are
the market step of C.3, with the active sets and $E$ fixed. If only $j$'s costs change,
$d\ln P=\varpi_j\hat c_j$, so $d\ln S_j=-(\sigma-1)\Lambda_j(1-\varpi_j)\hat c_j$ and,
for $h\neq j$, $d\ln S_h=(\sigma-1)\Lambda_h\varpi_j\hat c_j$. Prices are
$p_i=\mu(S_g)c_i$, so $d\ln p_i=\hat c_i+\varepsilon_\mu(S_g)\,d\ln S_g$, Averaging over
group $j$'s plants with weights $s_i/S_j$ gives $\rho_j\hat c_j$, and every rival's
price rises. Since
$(\sigma-1)\varepsilon_\mu\Lambda_j<1$ and $1-\varpi_j\in[0,1)$, $\rho_j\in(0,1]$.
Under quantity competition with $\eta=1$, $\varepsilon_\mu=S/(1-S)$, which gives the
stated $\Lambda_g$. \qed

\paragraph{D.2 Proof of Proposition~\ref{prop:ownership}.} With $m(S)=(1+\kappa
S)/\sigma$, $\pi^H/E=\sigma^{-1}\sum_g\theta_gS_g(1+\kappa S_g)
=\sigma^{-1}[\bar\theta+\kappa\sum_g\theta_gS_g^2]$, and
$\sum_g\theta_gS_g^2=\bar\theta\cdot\mathrm{HHI}+\mathrm{Cov}_S(\theta,S)$ by
definition. Total profit is the case $\theta_g\equiv1$,
$\sum_g\pi_g/E=(1+\kappa\,\mathrm{HHI})/\sigma$, using $\sum_gS_g=1$. Subtracting
$\bar\theta\sum_g\pi_g$ from $\pi^H$ leaves $(E\kappa/\sigma)\mathrm{Cov}_S(\theta,S)$.
The constant-markup benchmark is the case $\kappa=0$. \qed

\paragraph{D.3 Proof of Proposition~\ref{prop:response}.} By \eqref{eq:profit},
$d\pi_g=(E/\sigma)(1+2\kappa S_g)\,dS_g$. Let $Z=\sum_hS_h\Lambda_h$, so
$\varpi_j=S_j\Lambda_j/Z$. From D.1,
$dS_j=-(\sigma-1)S_j\Lambda_j(1-\varpi_j)\hat c_j=-(\sigma-1)\varpi_j\hat c_j\sum_{h\neq
j}S_h\Lambda_h$ and $dS_h=(\sigma-1)\varpi_jS_h\Lambda_h\hat c_j$ for $h\neq j$.
Substituting into $d\pi^H=\sum_g\theta_gd\pi_g$ and collecting terms gives
\eqref{eq:response}. \qed

\paragraph{D.4 Proof of Corollary~\ref{cor:entry}.} Monotonicity in $A$ is Step 2 of
B.3. For spending, $\Omega(K;A)=E\,\tilde\Omega(K;A)$ with $\tilde\Omega$ strictly
increasing in $K$, so for two rungs with $K_1<K_2$ the net gain
$E[\tilde\Omega(K_2;A)-\tilde\Omega(K_1;A)]-[\Phi_2-\Phi_1]$ is strictly increasing in
$E$. So as $E$ rises, no rung below the group's current choice can become optimal,
and the chosen strength does not fall. \qed

\subsection{Solution Algorithm}
\label{app:algo}

\paragraph{E.1 Markets and entry.} Given an aggregate $A$, each group's share solves
the scalar equation \eqref{eq:sharemap}, and the market solves \eqref{eq:clearing} by
safeguarded Newton steps on $\ln A$, which never leave a bracket that contains the
unique root. With entry, each group's best response at a given $A$ is found by scanning
its ladder of efficient sets, and $A$ is found by bisection.

\paragraph{E.2 Prices, incomes and wages.} Given wages and active sets, price indices
are obtained by iterating on \eqref{eq:priceindex} combined with \eqref{eq:cost}--\eqref{eq:delivered}, and market sizes
and incomes by one linear solve. Wages are updated by the tâtonnement of
Theorem~\ref{thm:wages}, with a step within the bound in condition (ii), followed by
Newton steps once close to the solution.

\paragraph{E.3 Active sets.} At the equilibrium wages for the current configuration,
the groups with the largest gains from changing their sets do so, and the procedure
repeats until no group wishes to move. Moving only the groups with the largest gains
changes the path of the algorithm, not its fixed points, and does not amount to an
order of entry. When the loop does not reach a configuration in which no group wishes
to move, it returns the configuration with the fewest such groups and reports how many
there are. Because the joint equilibrium is not known to be unique, we solve from
several starting wage vectors and compare the configurations reached.

\end{document}
